\documentclass[10pt,a4paper]{article}

\usepackage{graphicx}
\usepackage{graphics}
\usepackage{color}
\usepackage[utf8]{inputenc}
\usepackage{times}
\usepackage{amsmath}
\usepackage[T1]{fontenc}

\title{Time-dependent Schroedinger equation (finite difference approach)}
\date{}

\begin{document}
\maketitle
\small

We use the 2D  harmonic oscillator from the previous assignment 
\begin{equation}
V(x,y)=\frac{1}{2}m\left(\omega_x^2 x^2 + \omega_y^2 y^2\right).
\end{equation}

We assume a rectangular grid of size $N \times N$. The choice of $N$ -- as large as possible for a calculation that lasts not more than 2 minutes.
We use the atomic units with $m=m_e=1$,  $\hbar=1$. Lets take $\omega_x=1$ and $\omega_y=1.001$ (a change with respect to the last week).

We will solve the Schroedinger equation \begin{equation} i\hbar \frac{\partial \Psi}{\partial t}=H \Psi\end{equation}
for a time-independent potential. 

The wave function in the initial moment is given by \begin{equation}\Psi(t=0)=\frac{1}{\sqrt{3}}\left(\Psi_0+\Psi_1+i\Psi_2\right),\end{equation}
Where $\Psi_0, \Psi_1,\Psi_2$ are the Hamiltonian eigenstates -- the lowest-energy state, the first-, and the second-excited states, respectively.  



\section*{Crank-Nicolson scheme}
We use the Crank-Nicolson scheme for solution of the Schroedinger equation
\begin{equation}
\Psi(t+\Delta t)=\Psi(t)+\frac{\Delta t}{2 i\hbar}\left(H\Psi(t)+H\Psi(t+\Delta t)\right)
\end{equation}

Lets use an iterative scheme for the time-stepping
\begin{equation}
\Psi(t+\Delta t)^{\nu+1}=\Psi(t)+\frac{\Delta t}{2 i\hbar}\left(H\Psi(t)+H\Psi(t+\Delta t)^\nu\right).
\end{equation}
We can take $\Psi(t+\Delta t)^0=\Psi(t)$.

The choice of the time-step and the number of iterations depend on our choice of $N$.
For $\Delta x=\frac{1}{8}$ and $\Delta t=0.001$ the number of iterations 
necessary for convergence is $\simeq 5$.

\section*{Tasks}
\begin{itemize}
\item For our initial condition the motion of the packet should be periodic with period ${2\pi}$. 
\item Lets cover the time over 2 nominal periods of motion. 
\item 
Plot the norm and the average energy as a function of time. 
\item Plot the average $\langle x \rangle (t)$ and $\langle y\rangle (t)$. 
\item Is $\langle y (\langle x\rangle ) \rangle$ a closed loop? 
\item Can you make the wave packet turn in the other direction ?
\item Can you make it oscillate in $x$ or $y$ direction without a rotation?
\item What happens if $\Psi(0)=\Psi_i$ (the eigenstate is set as the initial state)?
\end{itemize}

\end{document}
